\section{一般测度的结构与积分}
\label{sec:measure-general-structure}

在实轴和欧氏空间上，“一块区域有多大”是由长度、面积、体积直接给出的。
换到一般对象的集合后，这种天然的大小不存在了，我们要自己把测量的规则补回来。
本节做四件准备工作：先说清哪些集合允许被测量、再说明怎样给这些集合配质量；然后引入开集、闭集等拓扑概念在测量中的位置；接着把熟悉的积分搬到一般测度上；最后处理两个空间的联合结构。
这些内容在实轴上都已经出现过，这里不重复证明，重点是看清“换了底集以后，哪些结论原封不动，哪些条件必须重新声明”。

\subsection{可测结构与基本赋权}

拿到一个由任意对象组成的集合后，测量要分两步走，而且两步互不相干。
第一步决定\emph{哪些集合看得见、可以测量}，这是可测结构；第二步决定\emph{看得见的集合各自有多大}，这是测度。
同一套可见集合上可以配完全不同的质量，同一个质量的想法在更粗或更细的可见范围下也要分别讨论。
本小节先建立这两层，再给出点质量、测度混合和零测例外的说法。

\subsubsection{可测范围与子$\sigma$代数}

\paragraph{可测空间与生成结构}

设$X$是我们关心的全部对象，例如若干类别、若干样本或若干状态。
要测量它的子集，先要圈定一个“允许测量的集合族”。
为什么不能直接宣布所有子集都可测？
在实轴上已经看到，一旦允许不可数次地分割、覆盖和取极限，就会出现无法自洽赋长度的集合；因此允许范围必须满足一定封闭性，但又不能大到没有限制。

我们沿用$\sigma$代数的要求：集合族$\mathcal A\subseteq\mathcal P(X)$包含全集$X$，并且对相对补集和可数并封闭。
取补和可数并还自动带来另外几种封闭：空集属于$\mathcal A$，可数交
\[
 \bigcap_nA_n=X\setminus\bigcup_n(X\setminus A_n)
\]
属于$\mathcal A$，两个可测集的差$A\setminus B=A\cap(X\setminus B)$也属于$\mathcal A$。
换句话说，凡是能用可测集合作可数多次集合运算得到的集合，仍然允许测量。
配备了这样一个集合族以后，二元组$(X,\mathcal A)$称为可测空间，$\mathcal A$中的集合称为可测集合。

实际问题里，我们往往只先点名少数几个必须能测量的集合，例如区间、矩形或某些类别，其余集合由它们“生成”。
对一族集合$\mathcal C$，包含它的$\sigma$代数可能有很多，其中最小的一个记为$\sigma(\mathcal C)$。
最小性可以这样理解：幂集$\mathcal P(X)$总是一个包含$\mathcal C$的$\sigma$代数，所有这种$\sigma$代数的交仍然包含$X$且对补集、可数并封闭，所以交出来的最小范围确实存在。
这种“只规定少数起点，再封闭出全部可测集”的做法非常省事：后面验证一个映射是否可测、两项测度是否相同，都只需在起点集合上检查，再由封闭性推广到整个$\sigma$代数。

用一个有限例子体会“可见”的含义。
在$X=\{1,2,3,4\}$上，幂集$\mathcal P(X)$允许区分每一个点；但由划分$\bigl\{\{1,2\},\{3,4\}\bigr\}$生成的$\sigma$代数只有
\[
 \bigl\{\varnothing,\{1,2\},\{3,4\},X\bigr\}.
\]
在这套范围里，“第$1$个点”单独是看不见的，能测量的只有前两个点、后两个点这两个分块。
可见集合的多少，决定了我们能在多大精度上讨论问题。

\paragraph{子$\sigma$代数与同一底集上的粗化}

有时底集不变，但我们主动“眯起眼睛”，只看其中一部分集合。
若$\mathcal G\subseteq\mathcal A$本身也是一个$\sigma$代数，就称它为\term{子σ代数}。
$(X,\mathcal G)$中的对象一个不少，仍是原来的$X$，只是可见的集合变少了。
这种操作叫粗化，它和“把底集缩小成子集$Y\subseteq X$”不同：缩小底集是扔掉一批对象，粗化只是降低对同一批对象的分辨能力。

仍以四点集合为例。
在$\mathcal A=\mathcal P(X)$下四个单点都能单独测量；换成
\[
 \mathcal G=\{\varnothing,\{1,2\},\{3,4\},X\},
\]
就只能称“前两点合起来多重”“后两点合起来多重”，再也无法分开第$1$点和第$2$点。
若原来在$\mathcal A$上有测度$\mu$，只在$\mathcal G$上读取它，记作$\mu|_{\mathcal G}$，这只是缩小赋权规则的定义域，并没有删除任何点。
后面讨论\term{条件积分}时，会专门研究“在只能看到$\mathcal G$的情况下，怎样保留各分块上的累计信息”；这里先把粗化这种结构本身建立起来。

\subsubsection{测度、混合与零测例外}

\paragraph{一般测度及有限性条件}

圈定可见集合以后，就要给每个集合配一个非负的“大小”。
设$\mu$把每个可测集映到$[0,+\infty]$。
要求它满足两条：空集的大小为$0$；把互不相交的可数个集合合并时，总大小等于各部分大小之和，即
\[
 \mu\!\left(\bigcup_{n=1}^{\infty}A_n\right)=\sum_{n=1}^{\infty}\mu(A_n).
\]
这样的$\mu$称为测度，$(X,\mathcal A,\mu)$称为测度空间。
第二条（可数可加）是测量的核心：它保证“分成几块分别称，再合起来称”，结果一致。
底集换成一般对象后，这个直觉与长度完全相同；由它还能立刻得到单调性$A\subseteq B\Rightarrow\mu(A)\leq\mu(B)$和可数次可加，证明与实轴上一样，不再重复。

后面的定理对“总质量有多大”有不同要求，这里先把三个常用名词一次分清，免得以后混淆。
若$\mu(X)<\infty$，称$\mu$为\term{有限测度}；进一步若$\mu(X)=1$，称$\mu$为\term{概率测度}，此时每个集合的质量都在$[0,1]$内。
还有一个更弱、却非常常用的条件叫$\sigma$有限，满足它的测度即\term{σ有限测度}：可以用可数个可测集$E_k$盖住$X$，使每一块的质量都有限。
三者关系是：概率测度一定有限，有限测度一定$\sigma$有限（取$E_1=X$），但反过来不成立。
自然数集上的计数测度给每个点质量$1$、全空间质量无穷，却能被单点集$\{1\},\ldots,\{k\}$有限地盖住；实轴上的勒贝格测度总质量无穷，但区间$[-k,k]$每一块都有限——它们都是$\sigma$有限而不有限。

集合列的连续性沿用前面的结论，但要留意有限性出现在哪里。
集合逐渐扩大时，$A_n\uparrow A$就有$\mu(A_n)\uparrow\mu(A)$，不要求任何有限性；集合逐渐缩小时，$A_n\downarrow A$要得到$\mu(A_n)\downarrow\mu(A)$，必须先有$\mu(A_1)<\infty$。
否则$(n,\infty)\downarrow\varnothing$，而每一项的勒贝格长度都是无穷，无法相减得到零。
请记住：$\sigma$有限不是所有运算的万能前提，密度表示、乘积构造和有界函数积分各自需要的条件，会在用到时分别说明。

\paragraph{加权点测度与\term{可数混合}}

最容易想象的测度，是把质量直接放在若干个点上。
固定一个点$x$（并假设$\{x\}$可测），定义
\[
 \delta_x(A)=\mathbf 1_A(x)=
 \begin{cases}1,&x\in A,\\0,&x\notin A,\end{cases}
\]
称为集中在$x$的点质量：集合包含$x$就得质量$1$，否则得$0$。
它对应“全部重量压在一个点上”的极端赋权。

再取一列点$x_i$和非负权重$w_i\geq0$，把若干点质量加起来：
\[
 \mu=\sum_i w_i\delta_{x_i},\qquad
 \mu(A)=\sum_i w_i\,\mathbf 1_A(x_i).
\]
因为每个点属于哪些不交集合是唯一确定的，两个非负求和可以交换次序，所以$\mu$仍然可数可加；同一位置上的权重自动相加。
权重总和有限时$\mu$是有限测度，总和为$1$时就是概率测度。
这里有个细节：定义本身不要求每个单点都可测，但只要具体计算中用到$\delta_x(\{x\})=1$，就必须先确认$\{x\}\in\mathcal A$。

同样的“加权相加”还可以用在测度上。
设$\mu_i$是同一空间上的非负测度，取非负系数$a_i$，令
\[
 \mu(A)=\sum_i a_i\mu_i(A),
\]
非负级数换序保证$\mu$仍是测度。
若各$\mu_i$都是概率测度且系数和为$1$，得到的就是可数混合：以比例$a_i$在第$i$种随机规律之间随机挑选一种。
一个直观例子是
\[
 \mu=\tfrac12\,\lambda|_{[0,1]}+\tfrac12\delta_0:
\]
一半重量像水一样均匀铺在$[0,1]$上按长度分配，另一半重量全部压在点$0$上。
这样的混合提醒我们：一个测度完全可以同时拥有“连续铺开”的部分和“集中成点”的部分。

\paragraph{零测例外与完备化}

测量时，质量为零的集合允许我们忽略其中点上的例外。
可测集$N$若满足$\mu(N)=0$，称为零测集；某性质$\mu$几乎处处成立，是指不满足性质的点全部能装进某个可测零测集。
但“可以忽略”永远是相对于具体测度而言的：在勒贝格测度下单点$\{0\}$质量为零，可以忽略；在$\delta_0$下它的质量却是$1$，恰好是唯一重要的点。
还要注意例外能够合并的规模：可数个零测集的并仍是零测集，所以可数多个“几乎处处成立”的结论可以共用一个例外范围；不可数个零测集的并则没有这种保证。

另一个相关的性质是完备性。
若每个可测零集的任意子集都仍然可测，就称这个测度空间完备。
不完备时，可能有一个集合包含在零测集里，它本身却不可测——我们知道它“应当质量为零”，却无法在规则内给它赋值。
补救办法是完备化：把这些缺失的零测子集、以及它们与原可测集的并全部补进可测范围，并按外面的零包络赋零。
补进去以后，原来集合的测量值一个都不变，只是可测范围成为满足要求的最小扩展。
实轴上的具体构造前一章已经做过，这里只需理解它的作用：有了完备化，在零测集上随意改动一个函数的值，得到的函数仍然可测；否则“几乎处处相同”的两个函数可能一个可测、一个不可测。
因此使用几乎处处等价的函数时，要么挑一个可测的代表来用，要么明确进入完备化空间，不能简单地以为“反正是零测，怎样都行”。

\subsection{测度与拓扑结构}

仅有可测结构时，我们知道一个集合能否测量，却说不清它是开是闭、是远是近。
后面比较测度时，有些指标（例如全变差）完全不需要位置信息，另一些指标（例如Wasserstein）则必须知道两点之间的距离。
本小节把这两类条件分开：先说明开集怎样进入可测结构，再定义“测度集中在哪里”的\term{支持}。

\subsubsection{Borel结构与标准Borel范围}

\paragraph{从开集生成Borel结构}

当$X$本身带有拓扑（从而能谈开集、闭集）时，最自然的可测范围是让所有开集都可测，再封闭出其余集合。
这个由全体开集生成的$\sigma$代数称为Borel结构：
\[
 \mathcal B(X)=\sigma(\{U\subseteq X:U\text{ 是开集}\}).
\]
于是开集、闭集以及对它们作可数并、可数交、取补得到的集合都可测。
在实轴上，开区间都是开集，这个定义与前一章“由开区间生成$\mathcal B(\mathbb R)$”完全一致。
在离散空间里每个单点本身就是开集，所以Borel结构等于幂集；实轴的Borel结构却不等于全部子集，不可测集依然存在。
请特别注意：选定Borel结构只决定了“哪些集合可测”，它本身\emph{没有指定任何测度}；同一个$\mathcal B(\mathbb R)$上既可以放长度，也可以放正态分布或一个点质量。

选定一个\term{Borel测度}后再做完备化，可测范围还可能扩大。
例如勒贝格可测集族$=$实轴Borel集族$+$所有勒贝格零集的子集，比Borel结构更大。
因此要分清三个层次：拓扑规定开集，Borel$\sigma$代数由开集生成，完备化还要依赖具体选了哪个测度。
后文说“在标准Borel空间上成立”时，指的是中间这一层；不能不加核对地把完备化后才出现的集合代进去。

\paragraph{\term{Polish空间}与\term{标准Borel空间}}

后面的距离比较和测度分解经常需要一个“足够规整、又足够常见”的舞台，为此引入两个条件。
一个拓扑空间若既能找到可数个点稠密地铺在其中（可分，因而开集可以由可数个基本邻域拼出），又能配上一个与拓扑相容的完备度量（柯西列的极限仍落在空间内），就称为Polish空间。
本章用到\term{Wasserstein距离}时，就直接在“给定度量完备且可分”的空间上工作。
这两个条件各管一件事，谁也代替不了谁，更不可能仅仅因为“存在Borel$\sigma$代数”就自动成立。

可测结构之间也能谈“本质相同”。
若双射$h:(X,\mathcal A)\to(Y,\mathcal B)$把可测集映成可测集、且可测集的原像也都可测，也就是每个$B\in\mathcal B$有$h^{-1}(B)\in\mathcal A$、每个$A\in\mathcal A$有$h(A)\in\mathcal B$，则称$h$为\term{可测同构}。
它表示两边的“可测拼图”能一一对应地拼回去，不要求保持距离或其他几何性质。
若一个可测空间能与某个Polish空间的Borel空间可测同构，就称它为标准Borel空间。
实数、欧氏空间、有限或可数离散空间，以及Polish空间的Borel子集都是常见例子。
这里只陈述这一范围并直接使用，不证明其分类定理。
关键直觉是：标准Borel保证“测度上的拼图足够规整”，后面的测度分解等结论可以用；但它不指定距离，也不保证你随手选的距离完备，几何比较仍要另给度量。

\subsubsection{支持与\term{全支撑}}

知道一个测度在哪些集合上有质量后，还想直观地说出“质量集中在空间的哪一片”。
支持就是这个“质量所在地”的精确说法，它还能解释连续函数的几乎处处相等等价于什么。

\paragraph{邻域定义的支持及其全质量}

对可分度量空间上的有限Borel测度$\mu$，定义
\[
 \operatorname{supp}\mu
  =\{x\in X:\text{每个包含 }x\text{ 的开邻域 }U\text{ 都有 }\mu(U)>0\}.
\]
换句话说，支持里的点“走到多近都躲不开质量”；支持外的点则存在一个开邻域完全没有质量。
正因为支持外的点各自被一个零质量开邻域包住，这些邻域的并仍是开集，所以支持是闭集。
支持必须由邻域定义，而不能简单说成“正质量单点的集合”——连续分布下每个单点质量都是零，支持却可以是整个区间。

支持还几乎承载了全部质量。
利用可分性给出的可数拓扑基$\{V_j\}$：对支持外的每个点，先取一个零质量开邻域，再从中选出一个包含该点的基开集。
这样支持外的全部点就被可数个零质量基开集盖住，由可数次可加性，
\[
 \mu(X\setminus\operatorname{supp}\mu)=0.
\]
这一步“先用可数基把任意覆盖变成可数覆盖”不能省略——直接把不可数个零测开集相加是不合法的。
作为对照：均匀分布在$[0,1]$上的测度，每个单点质量都是零，支持却是整个$[0,1]$；$\delta_0$每个邻域只有在含$0$时才有质量，支持就是$\{0\}$。
支持、正质量单点的集合、以及任意一个质量为全重的集合，是三个不同概念，不要混为一谈。

\paragraph{全支撑与连续函数的逐点辨认}

若支持扩到整个空间，$\operatorname{supp}\mu=X$，等价于每个非空开集都有正质量，称$\mu$具有全支撑。
这个条件有一个很实用的推论：在全支撑下，两个连续函数若几乎处处相等，就处处相等。
反证很直接。设连续函数$f,g$几乎处处相等，却在$x_0$处不同。
连续性保证差值$|f-g|$在$x_0$的某个开邻域$U$内仍然严格为正；全支撑又给$\mu(U)>0$。
于是$f\neq g$的地方有正质量，与几乎处处相等矛盾。
这说明：把连续函数看成“只差零测集”的等价类时，全支撑能帮我们找回逐点唯一的代表。

没有全支撑时结论就要打折扣。
在$\delta_0$下，$f(x)=x$与$g(x)=0$几乎处处相等（它们只在$x=0$以外不同，而那里没有质量），却显然不是同一个函数。
一般地，只要在支持上连续且相等，就足以推出几乎处处相等；但要想把结论恢复到整个空间的每一个点，才需要全支撑这类额外条件。

\subsection{一般积分的适用条件}

积分的构造思想在一般测度上并不改变：先在取值简单的函数上按集合质量加权，再用简单函数逼近一般函数。
本小节只回答三个问题：函数满足什么条件才能与可测结构配合、积分什么时候是有限实数、以及前一章的三条积分极限定理在一般空间还成不成立。

\subsubsection{\term{可测映射}与积分条件}

\paragraph{原像判据与可测复合}

函数或映射要参与测量，先要保证它不会把输出端的可测事件拉回输入端的不可测集合。
设$f:(X,\mathcal A)\to(Y,\mathcal B)$，若对每个$B\in\mathcal B$都有
\[
 f^{-1}(B)=\{x\in X:f(x)\in B\}\in\mathcal A,
\]
就称$f$可测。
这里的$f^{-1}(B)$只是“原像集合”的记号，完全不要求$f$有逆函数。

为什么只需在生成族上检查？因为取原像与集合运算可交换：
\[
 f^{-1}(Y\setminus B)=X\setminus f^{-1}(B),\qquad
 f^{-1}\!\left(\bigcup_nB_n\right)=\bigcup_nf^{-1}(B_n).
\]
只要生成$\mathcal B$的那些集合的原像都可测，反复取补和可数并以后，整个$\mathcal B$中集合的原像就都可测。
请注意这与后面的“测度唯一性”是两种推理：前者只靠原像保持集合运算，后者还要靠测度的可加性和有限性。

可测映射复合后仍可测，因为$(g\circ f)^{-1}(C)=f^{-1}(g^{-1}(C))$；连续映射在两边都取Borel结构时自然可测。
对实值函数$f:X\to\mathbb R$，判据还能进一步简化：半直线$\{t:t>a\}$随有理数$a$生成Borel结构，所以只需验证$\{f>a\}\in\mathcal A$对所有有理数$a$成立。
这就把抽象的“可测映射”接回了前一章熟悉的“可测函数”，也为后面计算推前测度和测度核铺平了道路。

\paragraph{非负积分与绝对可积条件}

对非负可测函数，积分仍按“分层称重”定义。
若简单函数$s=\sum_i a_i\mathbf 1_{A_i}$中各层$A_i$可测、两两不交，系数$a_i\geq0$有限，就令
\[
 \int_X s\,d\mu=\sum_i a_i\mu(A_i).
\]
这就是“函数在第$i$层取$a_i$，该层质量为$\mu(A_i)$，乘积相加”。
一般非负函数由简单函数从下逼近，积分取这些逼近值的上确界，允许结果为$+\infty$；二进制逼近和“积分与简单表示无关”的论证与实轴相同，不再重写。

允许函数取负值时，要拆成正负两部分：
\[
 f^+=\max\{f,0\},\qquad f^-=\max\{-f,0\},\qquad f=f^+-f^-.
\]
两部分各自作非负积分，只在它们不都为$+\infty$时才相减。
如果
\[
 \int_X|f|\,d\mu<\infty,
\]
称$f$绝对可积；这时两部分积分都有限，$\int f\,d\mu$才是一个有意义的有限实数，线性和三角不等式等性质也随之成立。
一个常用的充分条件是：在有限测度上，有界可测函数一定可积，因为
\[
 \int_X|f|\,d\mu\leq\|f\|_\infty\,\mu(X)<\infty.
\]
这套积分构造本身并不要求$\sigma$有限。
另外提醒：有正有负的是\emph{被积函数}；\ref{sec:measure-signed-total-variation}小节把测度本身允许取负值，得到的是另一个对象（有符号测度），两者不要混淆。

积分对测度的可数混合也是线性的。
若$\mu=\sum_i a_i\mu_i$且$f\geq0$可测，则
\[
 \int_X f\,d\mu=\sum_i a_i\int_X f\,d\mu_i.
\]
这个等式先对指示函数成立（就是测度定义），再对简单函数由有限可加成立，最后对一般非负函数由单调逼近成立。
对有符号$f$，只有在合成测度下绝对可积时才能这样拆开。
取点质量混合时，它退化成最直观的加权求和：
\[
 \int_X f\,d\!\left(\sum_i w_i\delta_{x_i}\right)=\sum_i w_i f(x_i).
\]
一般空间上的积分因此统一了“连续铺开”的累计和“在离散点上”的加权平均。

\subsubsection{极限操作的可测性与积分交换}

使用函数列时要把两件事分开：先问极限函数本身能不能测，再问极限符号能不能搬进积分号。
两件事各有各的条件，“每一项都可测、都可积”并不自动保证它们的任意汇总合法。

\paragraph{可数极限与不可数上确界的边界}

对可数个可测函数，阈值集合的可数并、可数交仍然可测，所以
\[
 \sup_n f_n,\quad\inf_n f_n,\quad\limsup_n f_n,\quad\liminf_n f_n
\]
都是可测函数；逐点极限存在时，极限函数也可测。
非负级数的部分和是一种特殊的可数极限，因此前面“可数混合的积分可以逐项求和”有合法的函数对象作保障。

但$\sigma$代数只对\emph{可数}次运算封闭，不可数个可测函数的逐点上确界未必可测。
实轴Borel结构上有个很说明问题的反例：取一个非Borel集合$N$，每个单点示性函数$\mathbf 1_{\{x\}}$都是可测的，但
\[
 \sup_{x\in N}\mathbf 1_{\{x\}}=\mathbf 1_N
\]
却是不可测函数。
这划出了一般规则的边界：不可数次优化或取极限要保证可测性，必须在具体场景中另加可分性、可测选择等条件，不能因为“每一项都可测”就默认其不可数汇总也可测。

\paragraph{三条积分极限规则的适用条件}

前一章的三条积分极限定理，其证明只用到简单函数逼近、测度的可数可加和可数极限可测，不依赖实轴的几何，因此在任意一般测度空间原样成立。
这里不重复证明，只把使用条件强调清楚。

若非负可测函数列单调上升$f_n\uparrow f$，单调收敛定理给出
\[
 \int f_n\,d\mu\uparrow\int f\,d\mu,
\]
两边一起取$+\infty$也允许。
对任意非负可测函数列，Fatou不等式
\[
 \int\liminf_n f_n\,d\mu\leq\liminf_n\int f_n\,d\mu
\]
始终成立，它只给“先取极限”一侧的下界，甚至不要求$f_n$真的有逐点极限。
支配收敛定理则要求三个条件同时具备：$f_n$几乎处处收敛到可测的$f$、存在\emph{同一个}非负可积函数$g$、且$|f_n|\leq g$几乎处处成立；这时积分可以取有限极限，并且
\[
 \int|f_n-f|\,d\mu\to0.
\]
证明时可以把可数个“不收敛、不被支配、取无穷值”的例外合并成一个零测集，在其外使用Fatou式论证，从而避免无穷相减。
请特别注意“同一个支配函数”这个要求不能由“每个$f_n$各自可积”替代。
本小节始终固定测度、只改变被积函数；至于“测度本身变化时积分怎样收敛”，要等到\ref{sec:measure-sequence-convergence}节讨论。

\subsection{乘积结构与测度唯一性}

后面比较测度时，经常先在“矩形”这类容易计算的集合上构造或验证一个测度，再断言它在整个空间上被唯一确定；讨论两个空间的配对时，也需要联合空间的可测结构。
本小节建立三件共同工具：简单集合上的相等怎样推广到全部集合、两个测度空间怎样相乘、以及更多坐标下的柱集接口。
要先强调一点：给定两个边缘测度，乘积测度只是联合测度的一种，耦合和测度核会给出许多别的联合测度。

\subsubsection{生成族上的测度唯一性}

\paragraph{有限测度的$\pi$系统唯一性}

在半直线、矩形这类简单集合上赋值很容易，但我们真正想确定的是整个$\sigma$代数上的测度。
自然的问题是：如果在一族简单集合上两项测度取值相同，能否断定它们处处相同？
答案在一定条件下是肯定的。
若集合族$\mathcal C$对有限交封闭（$C_1,C_2\in\mathcal C\Rightarrow C_1\cap C_2\in\mathcal C$），称它为$\pi$系统；半直线和可测矩形都满足这个性质。

\begin{theorem}[$\pi$系统上的测度唯一性]
\label{thm:measure-pi-lambda-uniqueness}
设$\mathcal C$是包含$X$且生成$\mathcal A$的$\pi$系统，$\mu,\nu$是$(X,\mathcal A)$上的有限测度。
若$\mu(C)=\nu(C)$对每个$C\in\mathcal C$成立，则$\mu=\nu$在整个$\mathcal A$上成立。
\end{theorem}

“包含$X$”这一条不能漏，它保证两项测度的总质量相等。
证明的核心思路值得了解，因为它在乘积和核的构造中会反复用到。
令
\[
 \mathcal D=\{A\in\mathcal A:\mu(A)=\nu(A)\}
\]
为所有“两项测度取值相同”的集合。
因为总质量有限且相等，$A\in\mathcal D$时补集也相等：
\[
 \mu(X\setminus A)=\mu(X)-\mu(A)=\nu(X)-\nu(A)=\nu(X\setminus A).
\]
又因为测度可数可加，两两不交集合并起来后仍然相等，加上$X\in\mathcal D$，$\mathcal D$就是一个对补集和不交可数并封闭的$\lambda$系统。
$\pi$—$\lambda$定理说，包含$\pi$系统$\mathcal C$的最小$\lambda$系统恰好是$\sigma(\mathcal C)=\mathcal A$，所以全部可测集都在$\mathcal D$中。
直观上讲：在有限总质量下，只要在“积木”上对齐，整体质量、补集和不交拼接就全部对齐了。

$\sigma$有限版本要用，但必须写清共同的有限覆盖：存在$E_n\in\mathcal C$、$E_n\uparrow X$，使得$\mu(E_n)=\nu(E_n)<\infty$。
做法是先在每个$E_n$内用有限版本比较截断测度$\mu(\cdot\cap E_n)$与$\nu(\cdot\cap E_n)$，再令$n\to\infty$拼出全空间。
只知道两项测度各自$\sigma$有限、却没有同一套有限盖子，并不能保证它们在生成族上相等。

\subsubsection{有限乘积与累次积分}

\paragraph{乘积测度与一般Tonelli/Fubini接口}

要同时描述两个空间上的配对，先规定联合空间里哪些集合可测。
设$(X,\mathcal A,\mu)$与$(Y,\mathcal B,\nu)$都是$\sigma$有限测度空间。
可测矩形$A\times B$生成乘积$\sigma$代数$\mathcal A\otimes\mathcal B$；在矩形上自然定义
\[
 (\mu\otimes\nu)(A\times B)=\mu(A)\nu(B),
\]
也就是“两块质量相乘”，其中任一因子质量为零时乘积记为零。
在$\sigma$有限条件下，这个矩形赋值能唯一扩展成$\mathcal A\otimes\mathcal B$上的乘积测度$\mu\otimes\nu$。
唯一性来自矩形族是$\pi$系统、且两个因子的有限覆盖$X_k,Y_l$给出联合空间的有限矩形覆盖$X_k\times Y_l$；存在性是矩形赋权的标准扩展，前一章在欧氏空间已经构造过，这里直接引用。
但请记住：乘积测度只是“两边独立地各算各的”这一种联合方式；同样的边缘完全可以对应许多非乘积的联合测度。

当$X,Y$都是可分度量空间时，拓扑Borel结构和乘积可测结构是一致的：
\[
 \mathcal B(X\times Y)=\mathcal B(X)\otimes\mathcal B(Y).
\]
一方面坐标投影连续，开集的原像可测，所以Borel矩形都在乘积拓扑的Borel结构中；另一方面，两个可数拓扑基的矩形构成乘积拓扑的可数基，任何开集都是可数个开矩形之并。
有了这个等式，后面在联合空间上谈Borel耦合、下半连续成本和弱极限时，用的都是同一套可测集合，不会出现“事件可测性对不上”的问题。

积分换序的定理直接沿用前一章，这里只确认一般空间下的条件。
对非负的乘积可测函数$f$，Tonelli定理保证：每个截面可测，内层积分作为另一坐标的函数可测，而且两种累次积分与联合积分相等，
\[
 \int_{X\times Y}f\,d(\mu\otimes\nu)
 =\int_X\!\left(\int_Y f(x,y)\,d\nu(y)\right)d\mu(x)
 =\int_Y\!\left(\int_X f(x,y)\,d\mu(x)\right)d\nu(y),
\]
即使积分值为$+\infty$也成立——这正是非负性带来的便利，不用担心正负无穷相消。
对有符号函数，则必须先用
\[
 \int_{X\times Y}|f|\,d(\mu\otimes\nu)<\infty
\]
保证绝对可积，Fubini定理才给出有限的换序结论；而且截面可积只对几乎处处的剩余坐标成立。
特别要避免一个常见错误：只看到某一种次序下的有符号累次积分存在，不足以换序，必须先确认联合绝对积分有限。

\paragraph{乘积完备化中的截面解释}

如果$f$只在乘积测度的完备化中可测，处理截面时要小心。
可以先选一个与$f$乘积几乎处处相同、但对原乘积$\sigma$代数可测的代表$\widetilde f$，对$\widetilde f$使用上述定理。
两者的差异集包含在某个原乘积可测零集$N$内；对$\mathbf 1_N$用Tonelli可知，$N$的截面在几乎处处的剩余坐标下可测且质量为零。
这一步给出了处理$f$截面所必需的零包络，不能直接假定差异集的截面在原来的因子结构中可测。
在这些“好”坐标上，$f$与$\widetilde f$的截面只在零截面的子集上不同；必要时还要进入因子测度的完备化，才能把这个子集当作可测零集。
所以不能宣称每个截面都合法，也不能由“两个因子分别完备”推出乘积$\sigma$代数本身已经完备。
把这层细微区别交代清楚以后，后面只在条件已经核对过的场合调用换序。

\subsubsection{多坐标与柱集的接口}

两个空间的乘积可以重复进行，得到任意有限个空间的乘积；但要描述无穷多个坐标（例如一条随时间展开的轨迹），得先说清可测范围和不同维数之间怎样保持一致。
这里只建立接口，不提前讲随机过程的扩张理论。

\paragraph{有限坐标信息与相容边缘}

在乘积空间$\prod_{i\in I}X_i$上，只约束有限多个坐标的集合称为柱集；全体柱集生成这里的乘积$\sigma$代数。
每个坐标投影都可测，只依赖有限个坐标的函数也仍可用原像判据检验可测性。
设想对每个有限指标集合$J$都指定一项联合测度$\mu_J$。
它们不能各自任意指定：从较大的有限坐标集投影回$J$以后，必须恰好恢复$\mu_J$，这就是相容性。
它保证“在更粗的有限观察下看到的边缘分布”，在不同维数的赋值中保持一致。

有限维相容是构造无穷坐标测度时必须通过的检查，但仅靠它不能在任意可测空间上断言扩张存在；还需要标准Borel这类规整条件，扩张定理、路径的正则性以及随机过程解释都留给后文的概率论。
到此为止，后面比较与映射所需的可测范围、一般赋权、积分条件和联合计算工具已经备齐；下一节起，我们固定同一个可测空间，开始比较上面不同测度之间的关系。
